\chapter{Introduction}
\label{ch:introduction}

An initial compromise can enable further modeled reachability under a specified network topology and policy. In this thesis, the \emph{blast radius} is the set of resources reached by the simulator after an initial compromise; it is a model output, not a measurement of real attacker movement~\cite{ou05mulval}. Conventional vulnerability prioritization often uses the CVSS score published by FIRST.org~\cite{cvss31}. CVSS describes vulnerability severity characteristics; it does not establish real-world exploit likelihood.

This thesis investigates a context-graph and state-transition approach, informed by attack-graph literature: represent network topology and CVE data, simulate propagation by Monte Carlo sampling, and use simulated mission-impact and blast-radius outcomes to rank defensive actions. The system is implemented in \mintinline{elixir}|Elixir|, with a visualization front-end in \mintinline{elixir}|Phoenix.LiveView|. It compares equal-action-count defense strategies with conventional CVSS-priority ordering; no full-study empirical comparison is reported in this thesis.

\begin{quote}
	\itshape A vulnerability score characterizes a flaw. A context graph defines the actions permitted by a specified model.
\end{quote}

The remainder of this chapter frames the goal of the thesis and outlines its structure. Section~\ref{ch:intro:goal} states the goal and scope. Section~\ref{ch:intro:questions} lists the research questions. Section~\ref{ch:intro:structure} describes the organization of the remaining chapters, and points forward to the architecture diagram in Chapter~\ref{ch:design} and the evaluation protocol in Chapter~\ref{ch:evaluation}.

Chapter~\ref{ch:theory} defines the context graph, state-transition process, and modeled blast-radius quantity used throughout the thesis.

\section{Goal and scope of the thesis}
\label{ch:intro:goal}

The goal of the thesis is to design and implement a small system, and to specify a reproducible evaluation protocol for it. The primary outcome is the simulated mission impact of an attack; blast radius is a secondary safety outcome. The current scope is intentionally limited to a single-pass offline analysis of one fixed stylized topology, policy, and attacker model. Reachability is canonical segment policy with operational flows derived in memory, and segmentation removes a policy rule. Each supported defensive action has unit cost, so the budget counts actions rather than deployable cost. The evaluation compares equal-action-count defense strategies with CVSS-priority ordering within an approved three-tier topology-scale study; current code supports the fixed baseline and generic topology sources, but not frozen controlled variants. The scope cannot establish real lateral-movement behavior, cost-aware segmentation, sensitivity to topology density or attacker profiles, or general scalability beyond measured scenarios. Runtime NVD ingestion is deferred; the baseline scenario may combine a reviewed static NVD subset with synthetic catalog entries. Online re-evaluation under topology change, live-CVE feed integration, and multi-tenant operation are out of scope.

\section{Research questions}
\label{ch:intro:questions}

The work is organized around one main research question and two sub-questions, which determine the evaluation protocol:

\begin{enumerate}
	\item \textbf{Main.} Within controlled stylized network topologies, how do equal-action-count defense strategies differ from CVSS prioritization in reducing simulated mission impact?
	\item \textbf{Sub-question 1.} How does a pre-attack feasibility constraint change the selected defense plans and the directly measured operational loss?
	\item \textbf{Sub-question 2.} How does controlled growth in topology size affect the simulation, strategy-selection, and total-evaluation runtime, and the mission-impact ranking of equal-budget defense strategies?
\end{enumerate}

\section{Structure of the thesis}
\label{ch:intro:structure}

The thesis is organized into seven chapters. Chapter~\ref{ch:literature_review} surveys prior work on attack graphs and vulnerability prioritization. Chapter~\ref{ch:theory} establishes the theoretical background and serves as a visual-elements reference. Chapter~\ref{ch:design} presents the design. Chapter~\ref{ch:implementation} describes the implementation. Chapter~\ref{ch:evaluation} specifies the evaluation protocol. Chapter~\ref{ch:conclusion} summarizes the contributions and outlines future work.
